\chapter{图形与直角}\label{ch:g3:shapes}

正方形、长方形、三角形、圆：这些图形都是老朋友。
本章用几何学家的眼光来看它们 --- 究竟
是什么让正方形成为正方形？ --- 并引入判断的工具：
三角尺及其直角。

\section{直角}

\begin{definition}[直角]\label{def:g3:shapes:rightangle}
\emph{直角}\index{直角}是纸张或三角尺
上的角。两条构成直角的直线
互相\emph{垂直}。图中，直角用
角上的一个小正方形标出。
\end{definition}

\begin{method}[检验直角]\label{met:g3:shapes:check}
\begin{enumerate}
  \item 将三角尺的直角顶点
        精确放在待检验的角顶点上；
  \item 把三角尺的一条边沿着角的一边滑动；
  \item 若角的另一边恰好沿三角尺的另一边
        重合 --- 则该角为直角；若张开更大或更小，
        则它不是直角。
\end{enumerate}
\end{method}

\begin{omfigure}
\begin{tikzpicture}[scale=0.85]
  \draw[thick] (0,0) -- (3,0);
  \draw[thick] (0,0) -- (0,2.2);
  \draw[thin] (0.32,0) -- (0.32,0.32) -- (0,0.32);
  \node[below, font=\small] at (1.5,-0.2) {一个直角};
  \begin{scope}[xshift=5.6cm]
  \draw[thick] (0,0) -- (3,0);
  \draw[thick] (0,0) -- (55:2.4);
  \node[below, font=\small] at (1.5,-0.2) {不是直角（太窄）};
  \end{scope}
  \begin{scope}[xshift=11.2cm]
  \draw[thick] (0,0) -- (3,0);
  \draw[thick] (0,0) -- (120:2.2);
  \node[below, font=\small] at (1.5,-0.2) {不是直角（太宽）};
  \end{scope}
\end{tikzpicture}

{\small 只有第一个角是直角 --- 而且只有小正方形标记，
或三角尺，才能确认它：仅凭肉眼
很容易看走眼。}
\end{omfigure}

\section{图形及其性质}

\begin{definition}[正方形、长方形、三角形]\label{def:g3:shapes:shapes}
\begin{itemize}
  \item \emph{长方形}\index{长方形}有 $4$ 条边，且有 $4$
        个直角；它的对边长度相等。
  \item \emph{正方形}\index{正方形}是 $4$ 条边
        长度都相等的长方形。
  \item \emph{三角形}\index{三角形}有 $3$ 条边；
        \emph{直角三角形}只有一个直角。
  \item \emph{圆}\index{圆}是用圆规画出的：
        圆上所有点到一个定点的距离都相同，这一距离称为
        \emph{半径}\index{半径}，该定点即\emph{圆心}。
\end{itemize}
\end{definition}

\begin{omfigure}
\begin{tikzpicture}[scale=0.72]
  % square, right-angle marks drawn inward at each corner
  \draw[thick] (0,0) rectangle (2,2);
  \foreach \x/\y/\dx/\dy in {0/0/1/1, 2/0/-1/1, 0/2/1/-1, 2/2/-1/-1}
    \draw[thin] (\x+0.22*\dx,\y) -- (\x+0.22*\dx,\y+0.22*\dy)
      -- (\x,\y+0.22*\dy);
  \draw (0.93,-0.07) -- (1.07,0.07);
  \draw (0.93,1.93) -- (1.07,2.07);
  \draw (-0.07,0.93) -- (0.07,1.07);
  \draw (1.93,0.93) -- (2.07,1.07);
  \node[below, font=\small] at (1,-0.35) {正方形};
  % rectangle
  \begin{scope}[xshift=3.8cm]
  \draw[thick] (0,0) rectangle (3,2);
  \foreach \x/\y/\dx/\dy in {0/0/1/1, 3/0/-1/1, 0/2/1/-1, 3/2/-1/-1}
    \draw[thin] (\x+0.22*\dx,\y) -- (\x+0.22*\dx,\y+0.22*\dy)
      -- (\x,\y+0.22*\dy);
  \node[below, font=\small] at (1.5,-0.35) {长方形};
  \end{scope}
  % right triangle
  \begin{scope}[xshift=8.6cm]
  \draw[thick] (0,0) -- (2.6,0) -- (0,2) -- cycle;
  \draw[thin] (0.28,0) -- (0.28,0.28) -- (0,0.28);
  \node[below, font=\small] at (1.3,-0.35) {直角三角形};
  \end{scope}
  % circle
  \begin{scope}[xshift=13.2cm]
  \draw[thick] (1,1) circle (1.1);
  \fill (1,1) circle (1.5pt);
  \draw[omDef, thick] (1,1) -- ++(45:1.1);
  \node[omDef, font=\small, above right] at (1.72,1.68) {半径};
  \node[below, font=\small] at (1,-0.35) {圆};
  \end{scope}
\end{tikzpicture}

{\small 图形一览，附带标记：小正方形表示直角，
小短线表示相等的边。}
\end{omfigure}

\begin{example}[标记为何重要]\label{ex:g3:shapes:marks}
画得稍微倾斜的图形，可能\emph{看上去}像正方形，
其实并不是。标记说明真相：$4$ 个直角标记 $+$ $4$ 个
短线标记 $=$ 真正的正方形。画图时要把标记画上；读图时，
只相信标记和给定的数据。
\end{example}

\section{在方格纸上作图}

\begin{method}[方格中的作图]\label{met:g3:shapes:grid}
方格纸自带直角：它的线互相垂直。
\begin{enumerate}
  \item 要画一个 $5$ 乘 $3$ 的长方形：选定一个交点，
        向右数 $5$ 格，向上数 $3$ 格，然后把图形
        沿方格线围成一圈；
  \item 要检验长度是否相等，就数一数方格；
  \item 斜着走的步子会画出倾斜的线：不断重复“向右 $1$ 格、向上 $1$ 格”，
        就能画出一条 $45^\circ$ 的直线。
\end{enumerate}
\end{method}

\begin{omfigure}
\begin{tikzpicture}[scale=0.5]
  \draw[gray!40, thin] (0,0) grid (13,5);
  \draw[very thick, omDef] (1,1) rectangle (6,4);
  \draw[very thick, omThm] (8,1) -- (12,1) -- (8,4) -- cycle;
  \node[omDef, font=\small, below] at (3.5,-0.15)
    {$5 \times 3$ 长方形};
  \node[omThm, font=\small, below] at (10,-0.15) {直角三角形};
\end{tikzpicture}

{\small 方格纸上的两种作图：方格线保证了
直角。}
\end{omfigure}

\begin{example}[作图程序]\label{ex:g3:shapes:program}
一个图形可以用一组指令来描述：
“1. 画一个正方形 $ABCD$，边长为 $4$ cm。2. 画一个圆，
圆心为 $A$、经过 $B$ 点。”按照程序逐行操作，
任何人都能重画出同一个图形 --- 而自己动手写程序，
会迫使你给每个点命名。和同伴试一试：一人写，
另一人来画。
\end{example}

\section{练习}

\begin{exercise}[$\star$]\label{exo:g3:shapes:1}
在教室里，不用量，找出三个直角和一个
不是直角的角。你怎样\emph{检验}每一个？
\end{exercise}

\begin{exercise}[$\star$]\label{exo:g3:shapes:2}
用直尺和三角尺画一个 $6$ cm 乘 $4$ cm 的长方形。标出
它的直角和相等的边。
\end{exercise}

\begin{exercise}[$\star$]\label{exo:g3:shapes:3}
画一个边长为 $5$ cm 的正方形。它有几个直角？
有几条相等的边？把它们都标出来。
\end{exercise}

\begin{exercise}[$\star$]\label{exo:g3:shapes:4}
在方格纸上画出：一个 $4 \times 4$ 的正方形；一个 $7 \times 2$
的长方形；一个直角边分别为 $3$ 格和 $5$ 格的直角三角形。
\end{exercise}

\begin{exercise}[$\star$]\label{exo:g3:shapes:5}
判断对错，并说明理由。
“每个正方形都是长方形。”\quad
“每个长方形都是正方形。”\quad
“一个三角形可以有两个直角。”（试着画一个！）
\end{exercise}

\begin{exercise}[$\star$]\label{exo:g3:shapes:6}
画一个半径为 $4$ cm 的圆。标出它的圆心 $O$、圆上的一点 $A$，
并画出半径 $[OA]$。$OA$ 有多长？如果 $B$
是圆上的另一点，$OB$ 有多长？
\end{exercise}

\begin{exercise}[$\star$]\label{exo:g3:shapes:7}
按照下面的程序操作：“1. 画一个长方形 $ABCD$，使 $AB = 6$ cm
且 $BC = 3$ cm。2. 标出点 $M$，它是 $[AB]$ 的中点。
3. 用线段把 $M$ 分别与 $C$、$D$ 连起来。”长方形被分成了
哪些图形？
\end{exercise}

\begin{exercise}[$\star$]\label{exo:g3:shapes:8}
给这些图形分类：哪些图形至少有一个直角？
正方形；（非正方形的）长方形；直角三角形；圆；
用粗线条画出的字母 L。
\end{exercise}

\begin{exercise}[$\star\star$]\label{exo:g3:shapes:9}
在方格纸上画一个倾斜的正方形，各边分别按“向右 $2$ 格、向上 $1$
格”的规律画出（转弯处同理）。数一数方格，
确认四条边的长度相等。
\end{exercise}

\begin{exercise}[$\star\star$]\label{exo:g3:shapes:10}
写出一个作图程序（要有编号的指令，并给点命名），
用于画出：边长为 $3$ cm 的正方形，立在一个长方形之上，
这个长方形为 $3$ cm $\times$ $2$ cm，好比房子的底层。用你的
程序请一位同学试一试。
\end{exercise}

\begin{exercise}[$\star\star$]\label{exo:g3:shapes:11}
在一个 $3 \times 3$ 的小方格网中，你能找出多少
个正方形？（请数一数 $1 \times 1$、$2 \times 2$ 和
$3 \times 3$ 的正方形。）
\end{exercise}
